\section{Hypothesis}

{\color{red} TODO: incoporate KG feedback from 10/29.  TODO: consider incorporating aspects of the hypothesis throughout appropriate sections.}

% From observing humans whipping a fixed-endpoint cable from left to right, we see that the trajectory of the cable's held end is a smooth curve that resembles a arc where the first half of the arc moves the held end up and the second half moves the held end down. Furthermore, we observe that even when humans whip the fixed-endpoint cable with a small reach, the middle segment of the cable often travels a much longer distance than the reach of the human arm. We can model this phenomenon using the following physical property: the speed of a wave $v$ along an unconstrained cable can be calculated as $v=\sqrt{T_c/\mu}$, where $T_c$ is the tension in cable and $\mu$ is the mass per unit length. When whipping an unconstrained cable homogeneous in mass, assuming minimal drag and friction, the velocity along the cable should be constant. In reality, when we are using a finite-length cable, the energy loss in drag and friction should be minimal. However, since we fix the endpoint of the cable, the endpoint's velocity is constrained to be 0, which will propagate the kinetic energy that should have been dissipated from the endpoint into the cable, making the middle segments of the cable travel faster, thus travelling a longer distance than the held end of the cable does in a given period of time. With this observation, we can whip the cable to long distances with a robot arm that has a fixed, limited reach. Additionally, at the apex of the motion, the held end experiences a shift in direction, from upwards to downwards, while still maintaining a non-zero horizontal velocity due to inertia. By changing the location of the apex of the whipping motion in space, we are also changing the kinetic energy and gravitational potential energy stored at the apex, thus changing the velocity induced along the cable during the descent part of the motion, varying the landing location of the cable. Thus, by changing the apex of the whipping motion, we can also change the landing location of the cable.
As one end of the cable is fixed and the robot can reset the cable to approximately the same starting state each time, we conjecture that for a reasonable range of motions, the trajectory of the cable can be repeatedly controlled by the trajectory of the robot gripper holding the other end of the cable. Fig.~\ref{fig:repeat} describes experiments on repeatability. The cable can always be reset to a relatively known initial state to the left of the obstacle. %To make the problem tractable and to avoid damage to the robot arm, 

We hypothesize that a minimum-jerk arcing motion defined by 3 points with max velocity at the apex will solve many vaulting problems. We instantiate the initial and final point for the robot at the far left and right of the robot's reachable workspace. The problem now is reduced to finding the apex (mid-point) for the arc and a minimum-jerk trajectory.
An illustration of this three-point arcing motion is shown in Fig.~\ref{fig:3point}. We propose the following procedure to dynamically manipulate the cable: given fixed starting and ending configurations of the robot arm holding the free end of a fixed-endpoint cable, generate a whipping motion by inferring the robot arm configuration at the apex of this whipping motion.
% Since the trajectory of the cable's held end when being whipped is almost arcing, and the robot arm has a fixed maximum reach, we take inspiration from parameterizing a arc of fixed $x$-intercepts and hypothesize that we can fix the starting and ending configurations of the robot arm and control the whipped cable by varying the robot configuration at the apex point of this semi-arcing motion. Besides validating the hypothesis in theory using the aforementioned energy view, we also validate the hypothesis in simulation by comparing the three-point arc model with different models. 

% We use simulation to evaluate the hypothesis on the 3-point dynamic motion, and to find a feasible base policy $\pi_\beta$ for data collection. Using simulation enables fast resetting to efficiently evaluate the qualitative behavior of multiple base policies. For simulation, we use Blender~\cite{blender}, because it has previously worked reasonably well for manipulating 1D~\cite{priya_2020} and 2D structures~\cite{descriptors_fabrics_2020}. Blender's keyframe feature allows us to specify robot configurations at certain animation frames, with the robot motion between keyframes being automatically interpolated. Additionally, Blender has an easy to use Python API, and its built-in rendering engine Eevee supports ray-tracing, making rendering fast on GPUs.

% To model cables in simulation, we use a Featherstone~\cite{featherstone1984robot} rigid-body simulator from BulletPhysics via Blender. Cables are implemented as a series of small capsule links connected by spring and torsional forces. We hand-tune the spring constants, torsional stiffness, and damping coefficients to make the simulated cable behave physically similar to a power cord. Specifically, in Blender, we use a torsional stiffness value of 1, a twist damping value of 0.9, a spring stiffness value of 40, and a spring damping value of 0.5.

% Using this setup, we compare the three-point arc model with different models. First, we tried regressing both apex point location and velocity explicitly, and found that regressing velocity did not improve the success rate of the vaulting task. We have also tried interpolating a straight-line motion between the points without making it arcing and found that the discontinuity in the middle causes the robot to come to a complete stop at the apex, thus reducing the energy stored for the cable to move fast enough. 


% We find the starting and ending configurations of the robot by moving the robot arm in free space to obtain the maximum left-to-right reach without violating safety constraints (e.g. hitting the camera mount). 

% We model the robot configuration at the apex as a 3D vector representing the three joint angles of the robot at the apex: $a_t = \langle \phi_{\rm base}, \phi_{\rm shoulder}, \phi_{\rm elbow} \rangle\in \mathbb{R}^3, \;\forall{t}\in [0,\dots T-1]$. In vaulting and knocking, $T=2$, and in weaving, $T=3$. 